\documentclass[10pt,a4paper]{amsart}
\usepackage[margin=19mm]{geometry}
\usepackage{amsmath,amssymb,mathtools}
\usepackage[scr=rsfs,cal=euler]{mathalfa}
\usepackage{xfrac}
\usepackage[unicode,hidelinks,bookmarks=false]{hyperref}
\newcommand{\ie}{i.e.\ }
\newcommand{\cf}{cf.\ }
\newcommand{\resp}{respectively\ }
\newcommand{\Subsection}{\subsection}
\providecommand{\nobreakdash}{\nobreak\hbox{-}}
\setlength{\emergencystretch}{2em}
\multlinegap0pt
\usepackage{amssymb}
\newtheorem{lem}{Lemma}
\newcommand{\C}{\mathbb{C}}
\renewcommand{\Im}{\operatorname{Im}}
\newcommand{\R}{\mathbb{R}}
\newcommand{\diag}{\operatorname{diag}}
\newcommand{\hj}[1]{h^{(1)}_{#1}}
\newcommand{\jj}[1]{j_{#1}}
\title{Determinant Formulas for Scattering Matrices of Schr\"odinger Operators with Finitely Many Concentric $\delta$-Shells}
\author{Masahiro Kaminaga}
\date{Source: arXiv:2603.24028v2 (CC BY 4.0)}
\begin{document}
\maketitle

\noindent\emph{Source and licence.} Determinant Formulas for Scattering Matrices of Schr\"odinger Operators with Finitely Many Concentric $\delta$-Shells. Version arXiv:2603.24028v2 (CC BY 4.0), distributed by arXiv under the Creative Commons Attribution 4.0 International licence, http://creativecommons.org/licenses/by/4.0/, as recorded in the arXiv licence metadata for this exact version and shown on the abstract page https://arxiv.org/abs/2603.24028v2. Full licence notice: \texttt{licenses/arxiv-2603.24028-CC-BY-4.0.txt}.\par\smallskip

\noindent\textit{Editorial scope.} Counted unit: the article's lem lem:partial-wave-construction (source0.tex lines 739--779). The article's complete original proof of it is delivered with it (source0.tex lines 781--1010, one \texttt{proof} environment of 230 lines). No mathematics is written, repaired or re-derived: the statement and the proof are the article's own lines, in the article's own notation, and no retained line is substituted or re-flowed.  Retained uncounted as the source-local context the proof needs: lines 396--415.  Nothing was omitted from any retained span, so the package carries no omission record.  No retained line contains a citation, so the wrapper carries no bibliography and no bibliography entry.  No retained line contains a figure environment, an image-inclusion command or drawing code, so no image is delivered and the package declares no asset dependency.  Editorial formatting only, in addition: an AMS \texttt{amsart} preamble that restates the article's own declarations and macros used by the retained lines, namely the environments \texttt{lem} on separate counters, together with the standard packages those lines need.  The retained environments print in this document as Lemma 1 in order of appearance, whereas the article prints the same items with the numbering of its own full document.  The complete native source of the article as distributed in the arXiv e-print for this exact version is bundled unchanged as \texttt{sources/197124/source0.tex}.
\begin{lem}\label{lem:mell}
Let $z\in\C\setminus[0,\infty)$ and set $k=\sqrt{z}$ with $\Im k>0$.
Then $m_\ell(z)$ is the $N\times N$ matrix whose entries are given by
\begin{equation}\label{eq:mell-entry}
(m_\ell(z))_{ij}
=
ik\,\jj{\ell}\!\bigl(k\min\{R_i,R_j\}\bigr)\,
\hj{\ell}\!\bigl(k\max\{R_i,R_j\}\bigr),
\qquad 1\le i,j\le N,
\end{equation}
where $\jj{\ell}=j_\ell$ is the spherical Bessel function and
$\hj{\ell}=h_\ell^{(1)}$ is the outgoing spherical Hankel function.
We also denote by $h_\ell^{(2)}$ the incoming spherical Hankel function.
In particular, $(m_\ell(z))_{ij}=(m_\ell(z))_{ji}$.
Moreover, the $\ell$--th partial--wave component of $K_N(z)$ is
\begin{equation}\label{eq:Kell}
K_\ell(z)=I_N+m_\ell(z)\Theta,
\end{equation}
where $I_N$ denotes the $N\times N$ identity matrix.
\end{lem}

\begin{lem}
\label{lem:partial-wave-construction}
Fix $\ell\ge0$, fix $m$ with $-\ell\le m\le\ell$, and let
$k>0$ satisfy
\[
\det K_\ell(k^2+i0)\neq0.
\]
Write
\[
K_\ell^+(k):=K_\ell(k^2+i0),
\qquad
b_\ell(k)
=
{}^t\bigl(\jj{\ell}(kR_1),\ldots,\jj{\ell}(kR_N)\bigr),
\]
and define
\[
c_\ell(k):=K_\ell^+(k)^{-1}b_\ell(k)\in\C^N.
\]
Then there exists an outgoing solution in the $(\ell,m)$ channel whose
incident part is
\[
\jj{\ell}(k|x|)Y_{\ell m}(\widehat x).
\]
For $|x|=r>R_N$, this solution has the form
\begin{equation}\label{eq:ext-form-lemma}
u_{\ell m}^+(x,k)
=
\frac12
\Bigl(
h_\ell^{(2)}(kr)+\sigma_\ell(k)\hj{\ell}(kr)
\Bigr)
Y_{\ell m}(\widehat x),
\end{equation}
where
\begin{equation}\label{eq:sigma-lemma}
\sigma_\ell(k)
=
1-2ik\,{}^t b_\ell(k)\Theta K_\ell^+(k)^{-1}b_\ell(k).
\end{equation}
\end{lem}

\begin{proof}
By assumption, the matrix
\[
K_\ell^+(k)=K_\ell(k^2+i0)
\]
is invertible, so $c_\ell(k)$ is well defined.

Let
\[
m_\ell^+(k):=m_\ell(k^2+i0),
\]
and let
\[
G_k^+(x,y):=\frac{e^{ik|x-y|}}{4\pi|x-y|}.
\]
For $j=1,\ldots,N$, define
\[
\Phi_j(x,k)
:=
\int_{S^2} G_k^+(x,R_j\omega')\,Y_{\ell m}(\omega')\,d\omega'.
\]
Equivalently,
\[
\Phi_j(x,k)
=
R_j^{-2}
\int_{S_j} G_k^+(x,y)\,Y_{\ell m}(y/R_j)\,d\sigma_j(y),
\]
where $d\sigma_j$ denotes the surface measure on $S_j$.

Thus $\Phi_j(\cdot,k)$ is smooth on $\R^3\setminus S_j$, satisfies
\[
(-\Delta-k^2)\Phi_j(\cdot,k)=0
\qquad \text{in }\R^3\setminus S_j,
\]
and is continuous across each sphere.

For $i\neq j$, the function $\Phi_j(\cdot,k)$ is smooth across $S_i$, hence
\[
\partial_r\Phi_j(R_i+0,\omega,k)-\partial_r\Phi_j(R_i-0,\omega,k)=0.
\]

For $i=j$, using the explicit formulas for $\Phi_j$ inside and outside $S_j$,
which will be derived below in \eqref{eq:Phi-exterior-lemma} and
\eqref{eq:Phi-interior-lemma}, we obtain, at $r=R_j$,
\begin{align*}
&\partial_r\Phi_j(R_j+0,\omega,k)-\partial_r\Phi_j(R_j-0,\omega,k)
\\
&\qquad
=
ik^2\Bigl(
j_\ell(kR_j)\bigl(h_\ell^{(1)}\bigr)'(kR_j)
-
h_\ell^{(1)}(kR_j)j_\ell'(kR_j)
\Bigr)Y_{\ell m}(\omega).
\end{align*}
Using the Wronskian identity
\[
j_\ell(t)\bigl(h_\ell^{(1)}\bigr)'(t)-j_\ell'(t)h_\ell^{(1)}(t)
=
\frac{i}{t^2},
\]
we obtain
\[
\partial_r\Phi_j(R_j+0,\omega,k)-\partial_r\Phi_j(R_j-0,\omega,k)
=
-R_j^{-2}Y_{\ell m}(\omega).
\]
Therefore
\begin{equation}\label{eq:jump-Phi-lemma}
\partial_r\Phi_j(R_i+0,\omega,k)-\partial_r\Phi_j(R_i-0,\omega,k)
=
-\delta_{ij}R_i^{-2}Y_{\ell m}(\omega).
\end{equation}

We next record two formulas for $\Phi_j$.

First, evaluating $\Phi_j$ on $S_i$ and using Lemma~\ref{lem:mell}, we obtain
\begin{equation}\label{eq:trace-Phi-lemma}
\Phi_j(R_i\omega,k)
=
(m_\ell^+(k))_{ij}Y_{\ell m}(\omega),
\qquad 1\le i,j\le N.
\end{equation}

Second, for $|x|=r>R_j$, the standard partial-wave expansion of the free Green
function yields
\[
G_k^+(x,R_j\omega')
=
ik\sum_{n=0}^\infty \sum_{q=-n}^n
j_n(kR_j)\,h_n^{(1)}(kr)\,
Y_{nq}(\widehat x)\overline{Y_{nq}(\omega')}.
\]
Multiplying by $Y_{\ell m}(\omega')$ and integrating over $S^2$, orthonormality of
the spherical harmonics gives
\begin{equation}\label{eq:Phi-exterior-lemma}
\Phi_j(x,k)
=
ik\,\jj{\ell}(kR_j)\hj{\ell}(kr)\,Y_{\ell m}(\widehat x),
\qquad |x|=r>R_j.
\end{equation}
Similarly, for $|x|=r<R_j$, the same partial-wave expansion gives
\begin{equation}\label{eq:Phi-interior-lemma}
\Phi_j(x,k)
=
ik\,\hj{\ell}(kR_j)\,\jj{\ell}(kr)\,Y_{\ell m}(\widehat x),
\qquad |x|=r<R_j.
\end{equation}
Hence $\Phi_j(\cdot,k)$ lies in the fixed $(\ell,m)$ channel on each annulus.

Now set
\begin{equation}\label{eq:scatt-sol-lemma}
u_{\ell m}^+(x,k)
=
\jj{\ell}(k|x|)Y_{\ell m}(\widehat x)
-
\sum_{j=1}^N \theta_j c_{\ell,j}(k)\,\Phi_j(x,k),
\end{equation}
where
\[
\theta_j=\alpha_jR_j^2,
\qquad j=1,\ldots,N,
\]
so that
\[
\Theta=\diag(\theta_1,\ldots,\theta_N).
\]
By \eqref{eq:Phi-exterior-lemma} and \eqref{eq:Phi-interior-lemma},
the function $u_{\ell m}^+(\cdot,k)$ also lies in the fixed $(\ell,m)$
channel on each annulus.

Since $\jj{\ell}(k|x|)Y_{\ell m}(\widehat x)$ is regular at the origin and each
$\Phi_j(\cdot,k)$ is smooth near the origin, the function
$u_{\ell m}^+(\cdot,k)$ is regular at $x=0$. Moreover,
\[
(-\Delta-k^2)u_{\ell m}^+(\cdot,k)=0
\qquad \text{in }\R^3\setminus\bigcup_{j=1}^N S_j.
\]

We verify the interface conditions. Evaluating
\eqref{eq:scatt-sol-lemma} on $S_i$ and using
\eqref{eq:trace-Phi-lemma}, we get
\[
u_{\ell m}^+(R_i\omega,k)
=
\Bigl(
\jj{\ell}(kR_i)-\sum_{j=1}^N (m_\ell^+(k)\Theta)_{ij}c_{\ell,j}(k)
\Bigr)Y_{\ell m}(\omega).
\]
Since
\[
K_\ell^+(k)c_\ell(k)=b_\ell(k),
\]
this becomes
\[
u_{\ell m}^+(R_i\omega,k)=c_{\ell,i}(k)Y_{\ell m}(\omega).
\]

On the other hand, \eqref{eq:jump-Phi-lemma} gives
\begin{align*}
&\partial_r u_{\ell m}^+(R_i+0,\omega,k)
-
\partial_r u_{\ell m}^+(R_i-0,\omega,k)
\\
&\qquad
=
-\sum_{j=1}^N \theta_j c_{\ell,j}(k)
\bigl(
\partial_r\Phi_j(R_i+0,\omega,k)
-
\partial_r\Phi_j(R_i-0,\omega,k)
\bigr)
\\
&\qquad
=
\theta_i c_{\ell,i}(k)R_i^{-2}Y_{\ell m}(\omega)
=
\alpha_i c_{\ell,i}(k)Y_{\ell m}(\omega)
=
\alpha_i\,u_{\ell m}^+(R_i\omega,k).
\end{align*}
Thus $u_{\ell m}^+(\cdot,k)$ satisfies the transmission conditions for the
$\delta$-shell interaction and solves
\[
(H-k^2)u_{\ell m}^+(\cdot,k)=0
\]
in the sense of distribution.

We next determine its exterior form. For $|x|=r>R_N$,
\eqref{eq:Phi-exterior-lemma} yields
\[
u_{\ell m}^+(x,k)
=
\Bigl(
\jj{\ell}(kr)
-
ik\,\hj{\ell}(kr)\,{}^t b_\ell(k)\Theta c_\ell(k)
\Bigr)Y_{\ell m}(\widehat x).
\]
Using
\[
\jj{\ell}(kr)=\frac12\bigl(\hj{\ell}(kr)+h_\ell^{(2)}(kr)\bigr),
\]
we obtain
\[
u_{\ell m}^+(x,k)
=
\frac12
\Bigl(
h_\ell^{(2)}(kr)+\sigma_\ell(k)\hj{\ell}(kr)
\Bigr)
Y_{\ell m}(\widehat x),
\qquad r>R_N,
\]
where
\[
\sigma_\ell(k)
=
1-2ik\,{}^t b_\ell(k)\Theta c_\ell(k).
\]
Since
\[
c_\ell(k)=K_\ell^+(k)^{-1}b_\ell(k),
\]
this is exactly \eqref{eq:sigma-lemma}.
By construction, the exterior part involves only $h_\ell^{(1)}(kr)$,
hence the solution is outgoing in the sense of the Sommerfeld
radiation condition.
\end{proof}

\end{document}
